\section{Square Digit Chains --- Problem 92}
\label{prob:92}
\noindent\textbf{Source:} \url{https://github.com/KO0I/projectEuler/blob/reorganize/solutions-by-language/solutions/c/square_digit_chains_92.c}

\noindent\textbf{Assessment:} Unimplemented.

\subsection*{Mathematical goal}
Starting from an integer \(n\), repeatedly replace it with the sum of the squares of its decimal digits. Count starting values below ten million whose chain reaches 89 rather than 1.

\subsection*{How the C code works}
The source includes a comment describing the digit-square transformation and example chains. Its \verb|main| contains no calculations, no loop over starting numbers, and no counter; it immediately returns zero.

\subsection*{Correctness and limitations}
This is a problem statement preserved as a C placeholder, not an algorithm. An implementation can precompute the destination of each possible digit-square sum and reuse that result across all starting values. There is no output from the current source that answers the problem; the mathematical count is \texttt{...1146}.


\subsection*{Correction to historical C}
{\color{red}\textbf{Correction:} The formerly empty C program now evaluates repeated decimal digit-square sums and counts exactly those starts below ten million that reach 89.}
